%%%% SELECT ONE OF THE FOLLOWING COMMANDS %%%%%%%%
%%% TEMPLATE FOR PROCEEDINGS TRACK %%%%
% \documentclass[mlmain]{jmlr}

%% TEMPLATE FOR Extended Abstract Track %%%%%%%
% Official limit: 4 pages, excluding references and appendices.
% (The Findings track is the one with no page limit.)
\documentclass[mlabstract]{jmlr}

%%%%%%%%%%%%%%%%%%%%%%%%%%%%%%%%%%%%%%%%%%%%%%%%%

% Watermark: required by the official review template; remove for camera-ready.
\usepackage[hpos=300px,vpos=70px]{draftwatermark}
\SetWatermarkText{\test}
\SetWatermarkScale{1}
\SetWatermarkAngle{0}

\usepackage{booktabs}

\jmlrvolume{}
\firstpageno{1}
\editors{NeurReps 2026 organisers}
\jmlryear{2026}
\jmlrworkshop{Symmetry and Geometry in Neural Representations}

\title[Curvature and local decodability]{Rank- and Scale-Conditioned Curvature Is Associated with Local Linear Decodability in Vision Representations}

% Double-blind: no author block.

\begin{document}
\maketitle

\begin{abstract}
Geometric summaries of neural activations depend on the rank and neighbourhood scale of the local chart used to compute them. We evaluate nested regularised quadratic charts of a ViT-B galaxy representation, define sphere-normal mean curvature $K_{H}^{\mathrm{cross}}$ from a split-half second fundamental form, and select chart ranks from held-out reconstruction rather than a single eigengap. At the frozen neighbourhood $k{=}2048$ ($n{=}512$ anchors), greater curvature is associated with poorer local out-of-fold linear decodability of $r$-band magnitude (\texttt{mag\_r\_desi\_local\_oof\_r2}) at the middle and upper evaluated ranks: controlled Spearman $\rho{=}-0.240$ at $d{=}16$ and $\rho{=}-0.233$ at $d{=}20$. Direct-error associations have the opposite sign (MAE $\rho{=}+0.251$, MSE $\rho{=}+0.227$), while local SST/target variance is essentially unassociated ($\rho{\approx}-0.025$). On $n{=}128$ hash-selected anchors the middle/upper direction remains negative at $k{=}1024$ and $k{=}1536$, but magnitudes attenuate substantially; $k{=}512$ fails curvature reliability and is not confirmatory. The association is therefore supported at the frozen scale and scale-dependent. Scope: one encoder, one astronomical dataset, one probe target, an exploratory rank sweep, and no independent replication.
\end{abstract}
\begin{keywords}
representation geometry, mean curvature, linear probes, vision transformers, chart rank
\end{keywords}

\section{Question and contributions}
\label{sec:intro}

Shared geometric structure is a working hypothesis for foundation-model activations~\citep{huh2024platonic,duraphe2025platonicuniverse}. We ask a narrower question about one encoder: after a nested quadratic chart is fit in a finite neighbourhood, is its sphere-normal mean curvature associated with how well a \emph{fixed} global linear probe~\citep{alain2017probes} reconstructs a physical scalar on that neighbourhood?

The probe is five-fold ridge regression from a ViT-B~\citep{dosovitskiy2021vit} embedding to $r$-band magnitude. The outcome at an anchor is the local out-of-fold coefficient of determination
\[
R^{2}_{\mathrm{local}}=1-\mathrm{SSE}_{\mathrm{OOF}}/\mathrm{SST}_{\mathrm{local}},
\]
denoted \texttt{mag\_r\_desi\_local\_oof\_r2}. Catalog magnitude \texttt{mag\_r\_desi\_catalog\_value} is a different vector and is never substituted. Curvature $K_{H}^{\mathrm{cross}}$ is the split-half inner product of mean-curvature vectors of a sphere-normal second fundamental form: finite-scale, extrinsic, and conditioned on chart rank and neighbourhood size.

\paragraph{Contributions.}
(i)~A method for evaluating local representation dimensionality as a predictive range rather than selecting one eigengap.
(ii)~A sphere-normal quadratic-chart curvature statistic that removes the forced radial curvature induced by activation normalisation.
(iii)~A controlled empirical association between rank-conditioned curvature and local OOF probe error at a frozen neighbourhood scale.
(iv)~An audit showing that the association is not an $R^{2}$-denominator artefact and is materially neighbourhood-scale dependent.

\section{Local quadratic geometry and chart rank}
\label{sec:geom}

At each of $512$ hash-stable anchors we take the $k{=}2048$ nearest neighbours on the unit sphere, remove the radial component, and compute a nested PCA frame $J$, following the local-quadratic charting tradition~\citep{donoho2003hessian}. A two-stage ridge quadratic is fit with split-half cross-evaluation: a tangential warp, then a second fundamental form $B^{S}$ projected onto the sphere-normal complement of $\mathrm{span}(x_{0},J)$. Packed off-diagonal coefficients store $2B_{ab}$. Mean curvature is $H_{a}=d^{-1}\sum_{i}B^{S}_{a,ii}$; the production scalar is $K_{H}^{\mathrm{cross}}=\langle H^{(A)},H^{(B)}\rangle$. Reliability $R_{H}$ is the split-half concordance of $H$. Identities appear in the appendix.

Held-out reconstruction (Fig.~\ref{fig:dim}) does not exhibit a single spectral jump. Pooled linear $R^{2}_{L}$ is approximately $0.80$ at $d{=}12$ and $0.85$ at $d{=}20$; quadratic NMSE decreases more slowly at higher ranks, while \emph{linear} reconstruction continues to improve at the right edge of the plotted family. We therefore treat $d{=}12$--$20$ as the \emph{evaluated chart range} spanning roughly $80$--$85\%$ explained variance, not as a complete geometrically adequate dimensional interval and not as an intrinsic dimension. A stable lower-dimensional core is compatible with further predictive gains at higher ranks. Chart positions used for scale mapping were fixed geometrically before inspecting curvature--probe correlations: lower $d{=}12$, middle $d{=}16$, upper $d{=}20$. Rank scanning inside the family is exploratory.

\begin{figure}[t]
\centering
\includegraphics[width=0.96\linewidth]{figures/figure1_dimension.pdf}
\caption{Chart rank is evaluated by held-out reconstruction, not a single eigengap.
Left: pooled linear and regularised quadratic NMSE versus chart rank ($k{=}2048$).
Right: pooled $R^{2}_{L}$ and its increment. Shading marks the \emph{evaluated chart range} $d{=}12$--$20$ ($\tau{\approx}0.80$ and $0.85$). Linear NMSE is still falling at $d{=}20$; no vertical line is drawn at a putative intrinsic dimension.}
\label{fig:dim}
\end{figure}

\section{Curvature and local decodability}
\label{sec:result}

\begin{figure}[t]
\centering
\includegraphics[width=0.92\linewidth]{figures/figure2_curvature_probe.pdf}
\caption{Finite-scale, rank-conditioned association between curvature and local probe error.
Left: controlled $\rho(K_{H}^{\mathrm{cross}},R^{2}_{\mathrm{local}})$ versus chart rank for the primary $n{=}512$, $k{=}2048$ analysis (line, simultaneous 95\% band). Markers show $n{=}128$ scale estimates; $k{=}512$ is marked unreliable. The top axis is held-out $R^{2}_{L}$. Shading: FWER-negative span $d{=}16$--$20$ in the primary analysis.
Right: percentile relationship between $K_{H}^{\mathrm{cross}}$ and OOF MSE at the predeclared representative rank $d{=}16$ (primary $n{=}512$, $k{=}2048$). Catalog magnitude and DESI labels are not shown.}
\label{fig:curv}
\end{figure}

\paragraph{Primary frozen analysis.} Aligning $K_{H}^{\mathrm{cross}}$ and \texttt{mag\_r\_desi\_local\_oof\_r2} by \texttt{sample\_id} on $n{=}512$ anchors at $k{=}2048$ reproduces the audited controlled Spearman correlations: $\rho_{12}{=}+0.143$, $\rho_{16}{=}-0.240$, $\rho_{20}{=}-0.233$ ($\rho_{20}{-}\rho_{12}{=}-0.376$). The corresponding raw (uncontrolled) correlation at $d{=}16$ is $\rho{=}-0.412$. The reconstructed $R^{2}_{\mathrm{local}}$ matches the frozen local-probe table to numerical identity. A typed schema rejects catalog magnitude as the primary target.

Controls are $\log$ $k$NN radius, local label variance, and evaluation count. Permutations ($B{=}10^{4}$) use rank-space Freedman--Lane for the controlled family; $p{=}0$ is reported as $p{<}1/(B{+}1)$. Familywise error is taken over ranks $d{=}12,\ldots,20$. The negative controlled associations at $d{=}16$ and $d{=}20$ have FWER $p{<}10^{-4}$. Paired bootstrap ($B{=}2000$) simultaneous 95\% intervals: $\rho\in[-0.323,-0.156]$ at $d{=}16$ and $\rho\in[-0.317,-0.147]$ at $d{=}20$. At $d{=}12$ the controlled association is positive ($\rho{=}+0.143$, FWER $p{=}0.013$). The rank that maximises $|\rho|$ is not treated as an intrinsic-dimension estimate.

\paragraph{Direct error versus denominator.} Greater curvature corresponds to \emph{higher} local prediction error and therefore \emph{lower} $R^{2}_{\mathrm{local}}$. After the same controls, at $d{=}16$, $k{=}2048$, $K_{H}^{\mathrm{cross}}$ is associated with OOF MAE ($\rho{=}+0.251$, 95\% CI $[0.154,0.333]$) and OOF MSE/SSE ($\rho{=}+0.227$, $[0.123,0.317]$), both FWER $p{<}10^{-4}$, and is essentially unassociated with SST / local target variance ($\rho{\approx}-0.025$, FWER $p{=}1.00$). NMSE equals $1-R^{2}_{\mathrm{local}}$ and is not an independent metric. Evaluating probe metrics on the outer neighbour half (indices $[k/2,k)$) while leaving $K_{H}^{\mathrm{cross}}$ unchanged yields $\rho_{16}{=}-0.246$ for $R^{2}_{\mathrm{local}}$ and $+0.238$ for MSE. A label shuffle at $d{=}16$ gives controlled $\rho{\approx}-0.059$.

\paragraph{Scale sensitivity.} The $n{=}512$, $k{=}2048$ result is the frozen primary analysis. Scale comparisons use $n{=}128$ hash-selected anchors and are not power-matched to the primary test. At the predeclared middle/upper ranks the direction remains negative at $k{=}1024$ and $k{=}1536$, but magnitudes attenuate substantially (controlled $\rho_{16}{=}-0.027$ and $-0.080$). The $k{=}2048$ common-$128$ estimate at $d{=}16$ is $\rho{=}-0.171$. Neighbourhood $k{=}512$ fails split-half curvature reliability (median $R_{H}{\approx}-0.19$) and is not interpreted as evidence of a contradictory scientific association. We do not apply post-hoc attenuation correction. Decision label: \texttt{claim\_supported\_but\_scale\_dependent}.
\FloatBarrier

\section{Limitations and conclusion}
\label{sec:limit}

$K_{H}^{\mathrm{cross}}$ is finite-scale extrinsic curvature of a fitted local chart of the activation cloud, not the intrinsic Riemannian curvature of an unknown manifold, and it is defined only at a chosen chart rank and neighbourhood size. The evaluated range $d{=}12$--$20$ does not identify a unique intrinsic dimension. The association is neighbourhood-scale sensitive. We study one encoder, one dataset, and one probe target. The rank sweep inside $d{=}8$--$20$ was not preregistered. There is no independent encoder or dataset replication. No causal interpretation is offered. DESI spectroscopic labels are excluded because sample-identity alignment was not established. Catalogue-magnitude correlations are not reported as discovery results.

\paragraph{Conclusion.}
At the frozen $k{=}2048$ scale, rank-conditioned sphere-normal mean curvature is associated with worse local OOF linear-probe performance at the middle and upper evaluated chart positions. The association survives direct-error and denominator checks, but its magnitude attenuates substantially at intermediate neighbourhood scales.

\clearpage
\bibliography{references}

\clearpage
\appendix
\section{Implementation notes}
\label{app:methods}

Definitions are taken from the code paths that produced the frozen artifacts; see also \texttt{METHODS\_FOR\_PAPER.md}.

Local coordinates: row-$\ell_2$ unit embeddings; $x_0$ normalised; radial component removed before nested PCA; $J\leftarrow\mathrm{orth}((I-x_0x_0^\top)J)$ (\texttt{sphere\_project\_basis}). Quadratic features $\phi_{ab}=u_au_b$ for $a\le b$; off-diagonal packed entries store $2B_{ab}$. $B^S$ columns are projected onto the complement of $\mathrm{span}(x_0,J)$. Production $K_H^{\mathrm{cross}}$ is the split-half inner product of $H$, not $\|H\|$. OOF probes are five-fold ridge ($\alpha{=}100$). Local $R^2$ uses uniform weights on finite neighbours:
\[
R^{2}_{\mathrm{local}}=1-\sum_i(y_i-\hat y_i)^2\Big/\sum_i(y_i-\bar y)^2.
\]
SSE is $\sum_i(y_i-\hat y_i)^2$; MSE is $\mathrm{SSE}/n_{\mathrm{eval}}$; MAE is $n_{\mathrm{eval}}^{-1}\sum_i|y_i-\hat y_i|$; SST is $\sum_i(y_i-\bar y)^2$. NMSE is $1-R^{2}_{\mathrm{local}}$ and is therefore not an independent estimand.

\section{Target identity and parity}
\label{app:parity}

The primary target is local OOF probe $R^{2}$ (\texttt{mag\_r\_desi\_local\_oof\_r2}). Catalog magnitude (\texttt{mag\_r\_desi\_catalog\_value}) is a distinct vector. A typed schema rejects ambiguous labels (\texttt{local\_r2}, \texttt{mag\_r\_desi}) and rejects a catalog vector that is identical or an affine rescaling of the probe-performance target. After \texttt{sample\_id} alignment, controlled Spearman correlations at $k{=}2048$, $n{=}512$ match the frozen audit to the reported rounding: $\rho_{12}{=}+0.143$, $\rho_{16}{=}-0.240$, $\rho_{20}{=}-0.233$. Reconstructed $R^{2}_{\mathrm{local}}$ matches the frozen geography table with maximum absolute difference $0$.

\section{Probe-metric associations at $k{=}2048$}
\label{app:metrics}

Table~\ref{tab:metrics} reports confirmatory controlled Spearman associations at the predeclared ranks (primary $n{=}512$). SST/variance remain near zero after controls; MAE/MSE reverse the $R^2$ sign, as required of an error association. Permutation $p{=}0$ is written $p{<}10^{-4}$ ($B{=}10^{4}$). Intervals are paired-bootstrap 95\% CIs; FWER is over $d{=}12,\ldots,20$.

\begin{table}[h]
\caption{Controlled $\rho(K_H^{\mathrm{cross}},\,\cdot)$ on $n{=}512$ anchors, $k{=}2048$.}
\label{tab:metrics}
\centering
\footnotesize
\begin{tabular}{lrrr}
\toprule
Target at $d{=}16$ & $\rho$ & 95\% CI & FWER $p$ \\
\midrule
\texttt{mag\_r\_desi\_local\_oof\_r2} & $-0.240$ & $[-0.323,-0.156]$ & ${<}10^{-4}$ \\
OOF MAE & $+0.251$ & $[0.154,0.333]$ & ${<}10^{-4}$ \\
OOF MSE / SSE & $+0.227$ & $[0.123,0.317]$ & ${<}10^{-4}$ \\
Local SST / target var.\ & $-0.025$ & --- & $1.00$ \\
NMSE $=1-R^{2}_{\mathrm{local}}$ & $+0.240$ & --- & --- \\
\midrule
$R^{2}_{\mathrm{local}}$ at $d{=}12$ & $+0.143$ & $[0.051,0.226]$ & $0.013$ \\
$R^{2}_{\mathrm{local}}$ at $d{=}20$ & $-0.233$ & $[-0.317,-0.147]$ & ${<}10^{-4}$ \\
\bottomrule
\end{tabular}
\end{table}

Raw (uncontrolled) Spearman at $d{=}16$ for $R^{2}_{\mathrm{local}}$ is $\rho{=}-0.412$. Outer-half neighbour probe metrics (indices $[k/2,k)$, $K_H^{\mathrm{cross}}$ unchanged) at $d{=}16$: $R^{2}_{\mathrm{local}}$ $\rho{=}-0.246$, MSE $\rho{=}+0.238$. Label shuffle at $d{=}16$: controlled $\rho{\approx}-0.059$.

\section{Scale subset and reliability}
\label{app:scale}

Chart positions were fixed from held-out variance ($\tau{=}0.80,0.85$ and the midpoint), not by maximising $|\rho|$. Table~\ref{tab:scale} uses the rank-sweep $128$-anchor refits. Signs at $d{=}16,20$ stay negative at $k{=}1024$ and $1536$, but magnitudes are ${<}50\%$ of the primary $k{=}2048$, $n{=}512$ value. Split-half reliability $R_{H}$ fails at $k{=}512$ (median $\approx-0.19$) and recovers to $\approx0.24$ at $k{=}1024$ and $\approx0.41$ at $k{=}1536$. The $k{=}512$ row is reported for completeness and is not confirmatory.

\begin{table}[h]
\caption{Controlled $\rho(K_H^{\mathrm{cross}},R^{2}_{\mathrm{local}})$ at predeclared ranks. Primary $k{=}2048$ uses $n{=}512$; other rows use $n{=}128$.}
\label{tab:scale}
\centering
\footnotesize
\begin{tabular}{lrrr}
\toprule
$k$ & $d{=}12$ & $d{=}16$ & $d{=}20$ \\
\midrule
$2048$ ($n{=}512$) & $+0.143$ & $-0.240$ & $-0.233$ \\
$512$ (unreliable $R_H$) & $-0.270$ & $-0.198$ & $-0.253$ \\
$1024$ & $+0.040$ & $-0.027$ & $-0.072$ \\
$1536$ & $+0.061$ & $-0.080$ & $-0.093$ \\
$2048$ common $128$ & $+0.148$ & $-0.171$ & $-0.193$ \\
\bottomrule
\end{tabular}
\end{table}

\section{Excluded analyses}
\label{app:excluded}

DESI spectroscopic labels are omitted because embedding-table identity alignment is unproven; equal row counts are not a join. Catalog-magnitude associations are retained only as a negative control against silent substitution and are not a discovery result. Adaptive-run physics labels other than the frozen probe target are not used for the claim. An earlier catalog-label substitution motivated strict target typing; it is not part of the scientific narrative of the main text.


\end{document}
